\chapter{附录A：高空间假阳性审计}

\section{控股平台与投资收益}

吉林敖东、辽宁成大等公司在部分模型中可能出现极高目标空间，来源往往是把投资收益或持股价值折价修复，与高成长经营PE混合。投资收益受被投企业利润、市场价格与会计处理影响，不具有可重复的单位经济。正确方法是资产净值折价、分红能力与持股价值敏感性，而不是25--45倍成长PE。

\section{一次性或低基数利润}

九安医疗等标的在特殊年份会产生巨大现金与利润，若用单期利润年化，目标价可以轻易翻倍。研究必须区分持续经营利润、一次性项目、资产处置、政府补助与投资收益。没有可重复业务桥，就只能做资产与现金价值分析。

\section{周期峰值}

江波龙是本轮最典型的峰值分母风险。Q1毛利率55.53\%，H1利润预告92--110亿元，但现金流显著为负，券商2027利润增长又很低。若存储价格回落或补库成本上升，当前看似便宜的2026PE会迅速失真。

锂盐标的也有类似问题。锂价、低价库存、自供矿与套保可能在某一季度共振；只有持续出货、正常化吨利、现金与资产负债表一起改善，才能给估值信用。

\section{拥挤成长}

生益科技与深南电路作为AI硬件对照，基本面可能保持强劲，但高价与高预期会显著抬高翻倍门槛。一个好行业不自动产生一倍空间。若市场已经完整资本化未来数年的增长，任何轻微低于预期都可能引发倍数收缩。

\begin{exhibitbox}[Exhibit A1：假阳性类型与修复方法]
\small
\begin{tabularx}{\exhibitboxwidth}{L{3cm} L{3cm} X X}
\toprule
假阳性 & 典型标的 & 错误方法 & 修复方法 \\
\midrule
投资收益当成长EPS & 吉林敖东、辽宁成大 & 高成长PE & NAV折价、分红与资产价值 \\
一次性利润年化 & 九安医疗 & 单期利润乘倍数 & 正常化经营利润与现金价值 \\
周期峰值永续 & 江波龙、锂盐公司 & 峰值EPS乘高PE & 单位经济、现金流、周期中枢 \\
热门主题忽略价格 & 生益科技、深南电路 & 景气即翻倍 & 当前隐含增长与估值拥挤 \\
关系/认证当订单 & 多数成长公司 & 客户名录即收入 & 公司订单、转量、回款与利润 \\
\bottomrule
\end{tabularx}
\end{exhibitbox}

\clearpage
